%% 自動生成（gen.py）
\bktocchap{6-12}{減衰振動を式で解く：特性方程式}{2次方程式の判別式が、ゆれ方を決める}{4}{5}{5}{420}{C}
\bktocchap{6-13}{ゆらし続けると：強制振動と共振}{小さな力でも、タイミングが合うと大きくゆれる}{4}{5}{5}{458}{C}
\bktocchap{6-14}{微分方程式がいくつも：連立微分方程式}{つながった2つの水そうの水位は、どう変わる？}{5}{5}{5}{492}{C}
\bktocchap{6-15}{まとめ：物理は「微分方程式＋はじめの状態」で書ける}{ばねも電気も熱も、同じ形の式で書ける}{1}{2}{3}{534}{C}
